\chapter{Mecánica de fractura y tolerancia al daño}\label{ch:03}
\objetivos{Relacionar energía y campo de punta; diferenciar $K_I$ de $K_{IC}$; calcular un tamaño crítico con una definición explícita de grieta y comprobar las restricciones del modelo.}
\section{Por qué no basta la resistencia nominal}
Una grieta modifica el campo de tensiones y concentra deformación en su punta. La tensión nominal media puede ser inferior al límite elástico y, sin embargo, existir propagación inestable. La mecánica de fractura introduce el tamaño y la geometría del defecto en el análisis. Su objeto no es afirmar que toda indicación sea crítica, sino cuantificar qué combinaciones de carga, defecto y resistencia a propagación son compatibles con el modelo elegido.

En fractura dúctil puede haber nucleación, crecimiento y coalescencia de microcavidades. En fractura frágil la deformación plástica macroscópica puede ser pequeña. Los términos dúctil y frágil describen comportamiento bajo unas condiciones: temperatura, espesor, velocidad de carga y restricciones pueden cambiarlo. La fractografía aporta indicios del mecanismo, pero exige relacionarlos con la historia y no confundir una superficie final de sobrecarga con el origen del daño.
\section{Modos de carga y convención de longitud}
El modo I abre las caras de la grieta; el modo II produce deslizamiento en el plano; el modo III, desgarro fuera de él. Un componente puede presentar modo mixto. Las fórmulas siguientes se limitan al modo I y a mecánica de fractura elástica lineal (MFEL).

Se define $a$ como \textbf{semilongitud} de una grieta central pasante de longitud total $2a$ en una placa idealizada muy ancha. Para una grieta de borde, $a$ suele representar su longitud desde el borde. La fórmula no resuelve esta ambigüedad por el estudiante: cada enunciado debe fijarla.
\ilustra{modos-grieta}{Esquemas de los modos I, II y III y definición de $a$ para una grieta central de longitud total $2a$.}
\section{Criterio energético}
La propagación crea superficie y modifica la energía potencial del sistema. La tasa de liberación de energía es $G=-\partial\Pi/\partial A_c$, con una convención consistente para el área de grieta creada. En un sólido idealmente frágil, el criterio de Griffith iguala energía liberada y energía requerida para crear superficies. En un material real, la resistencia incluye disipación y puede depender de la extensión de grieta; la curva R representa esa evolución.

Para modo I elástico lineal e isotropía,
\begin{equation}G=\frac{K_I^2}{E'},\qquad E'=\begin{cases}E&\text{tensión plana},\\ E/(1-\nu^2)&\text{deformación plana}.\end{cases}\label{eq:gk}\end{equation}
$G$ se mide en $\mathrm{J/m^2}$; $K_I$, en $\mathrm{Pa\sqrt m}$. La equivalencia energética no permite ignorar plasticidad extensa ni anisotropía. La expresión de $E'$ debe corresponder al estado mecánico adoptado.
\section{Intensidad aplicada y resistencia del material}
En un campo singular idealizado de punta, las tensiones elásticas escalan como $K_I/\sqrt{2\pi r}$ multiplicado por funciones angulares. La singularidad no afirma que un material soporte tensión infinita: cerca de la punta existe una zona de proceso que la idealización no resuelve.
\begin{equation}K_I=Y\sigma\sqrt{\pi a}.\label{eq:ki}\end{equation}
El factor $Y$ representa geometría, posición y carga. Puede depender de $a/W$ y no ser constante durante el crecimiento. La tensión $\sigma$ debe estar definida de forma coherente con la solución geométrica, por ejemplo sobre área bruta o neta. Para una placa ideal infinita con grieta central y tracción remota uniforme se adopta $Y=1$.

\textbf{$K_I$ es una solicitación calculada. $K_{IC}$ es una tenacidad crítica de modo I bajo condiciones de deformación plana y validez del ensayo.} El subíndice no es una distinción tipográfica prescindible. Una pieza delgada puede tener otra resistencia aparente $K_c$ dependiente del espesor; no se le atribuye automáticamente el valor $K_{IC}$ de una ficha.
\caja{Criterio de evaluación dentro de MFEL}{La comparación $K_I<K_{IC}$ solo informa sobre inestabilidad de modo I cuando son aplicables el campo elástico, la tenacidad y las condiciones de medida. No demuestra ausencia de crecimiento por fatiga, corrosión bajo tensión o fluencia.}
\section{Tamaño crítico, carga crítica y sensibilidad}
Si $Y$ y la tensión remota se mantienen constantes, imponer $K_I=K_{IC}$ conduce a
\begin{equation}a_c=\frac{1}{\pi}\left(\frac{K_{IC}}{Y\sigma}\right)^2,\qquad\sigma_c=\frac{K_{IC}}{Y\sqrt{\pi a}}.\label{eq:ac}\end{equation}
La derivación muestra que duplicar $\sigma$ divide $a_c$ entre cuatro, mientras que duplicar $K_{IC}$ multiplica $a_c$ por cuatro bajo las mismas hipótesis. Si $Y=Y(a/W)$, se resuelve $Y(a/W)\sigma\sqrt{\pi a}=K_{IC}$ de forma implícita y se comprueba que la raíz está dentro del dominio de la solución geométrica.
\caso{M03. No confundir $a$ con $2a$}{Una placa ideal muy ancha tiene una grieta central pasante de longitud total $4.0\,\mathrm{mm}$, $\sigma=180\,\mathrm{MPa}$, $Y=1$ y una tenacidad supuesta aplicable de $32\,\mathrm{MPa\sqrt m}$. Se usa $a=0.0020\,\mathrm m$. Así, $K_I=180\sqrt{\pi(0.0020)}=14.27\,\mathrm{MPa\sqrt m}$. El tamaño crítico es $a_c=(32/180)^2/\pi=0.01006\,\mathrm m$, es decir, longitud total crítica $2a_c=20.12\,\mathrm{mm}$. La sustitución $180\sqrt{\pi a_c}=32$ cierra el criterio. El resultado no autoriza mantener en servicio el defecto: faltan cargas variables, espesor, entorno y reglas de aceptación.
\par\textbf{Desarrollo para comparar tu intento.} El dato 4 mm es longitud total: la fórmula $K_I=Y\sigma\sqrt{\pi a}$ usa semilongitud. Para el tamaño crítico imponer $K_I=K_{IC}$, dividir por $Y\sigma$, elevar al cuadrado y dividir por $\pi$.

\begin{enumerate}

\item Convertir longitud total y dividir entre dos.

\[\fitmath{a=\frac{0{,}004}{2}\approx 0{,}002\,\mathrm{m}}\]

\item Evaluar intensidad con tensión en MPa y longitud en m.

\[\fitmath{\mathrm{KI}=180\cdot \sqrt{\pi\cdot 0{,}002}\approx 14{,}26798\,\mathrm{MPa\,m^{.5}}}\]

\item Despejar semilongitud crítica.

\[\fitmath{\mathrm{ac}=\frac{{\left(\frac{32}{180}\right)}^{2}}{\pi}\approx 0{,}01006016\,\mathrm{m}}\]

\item Volver a longitud total para comparar con el dato inicial.

\[\fitmath{\mathrm{total}=\frac{2\cdot {\left(\frac{32}{180}\right)}^{2}}{\pi}\approx 0{,}02012033\,\mathrm{m}}\]

\end{enumerate}

}
\section{Zona plástica y validez}
Una estimación de Irwin de la zona plástica de punta es $r_p\simeq(K_I/\sigma_y)^2/(2\pi)$ en tensión plana y, con mayor restricción, $r_p\simeq(K_I/\sigma_y)^2/(6\pi)$ en deformación plana. Son estimaciones, no límites exactos de una zona de geometría universal. La condición de pequeña plastificación exige que esa región sea pequeña respecto al tamaño de grieta, ligamento y dimensiones relevantes.

En ensayos normalizados se emplean condiciones dimensionales de validez. Una expresión de cribado habitual es
\begin{equation}B,\ a,\ W-a\ \gtrsim\ 2.5\left(\frac{K_{IC}}{\sigma_y}\right)^2.\label{eq:espesor}\end{equation}
No debe usarse como una certificación automática de un ensayo: la norma aplicable contiene otras condiciones, geometrías y criterios de carga. Aquí sirve para entender por qué tenacidad alta y límite elástico bajo pueden requerir una probeta grande.

Con $K=32\,\mathrm{MPa\sqrt m}$ y $\sigma_y=400\,\mathrm{MPa}$, la escala derecha vale $16\,\mathrm{mm}$. Una chapa de $3\,\mathrm{mm}$ no satisface ese cribado de espesor. Es incorrecto resolver el problema imponiendo deformación plana y declarar después que se ha comprobado solo porque el enunciado lo pedía.
\section{Cuando la elasticidad lineal resulta insuficiente}
Si la plasticidad no es pequeña pueden necesitarse parámetros elasto-plásticos como $J$ o CTOD. Para una respuesta elástica lineal, $J$ coincide con $G$ bajo las hipótesis correspondientes, pero ello no convierte $J$ y $K$ en símbolos intercambiables en cualquier situación. Una curva de resistencia al desgarro permite estudiar crecimiento estable y la interacción con la rigidez del sistema de carga.

La tolerancia al daño combina tamaño inicial supuesto, crecimiento, detectabilidad y tamaño final admisible. No utiliza el tamaño mínimo nominal de detección como garantía determinista: se consideran probabilidad de detección, accesibilidad, orientación y dispersión. El intervalo de inspección depende del tiempo de crecimiento entre estados justificadamente definidos y de márgenes adicionales.
\ilustra{fractura-critica}{Intensidad aplicada frente a semilongitud de grieta para el ejemplo M03. La intersección ilustra el criterio idealizado, no un límite de mantenimiento real.}

\section{La competencia entre crecimiento de grieta y plasticidad}
La intensidad elástica singular describe un campo intermedio: suficientemente lejos del proceso microscópico de la punta y suficientemente cerca para que domine el término de intensidad. Una estimación de escala de zona plástica en tensión plana es
\begin{equation}r_p\approx\frac{1}{2\pi}\left(\frac{K_I}{\sigma_y}\right)^2.\end{equation}
En una idealización de deformación plana se utiliza un factor diferente, frecuentemente $1/(6\pi)$. Son estimaciones de orden de magnitud. No sustituyen las condiciones completas de un ensayo ni permiten decidir por sí solas que la mecánica lineal es válida. Deben compararse con grieta, ligamento, espesor y gradientes.

Si la plasticidad deja de ser pequeña, energía liberada y deformación plástica se acoplan. Parámetros como la integral $J$ y la apertura de punta pueden ser más adecuados, bajo sus propias condiciones. Una curva de resistencia al crecimiento puede mostrar que la resistencia no permanece constante durante extensión estable. Un valor de $K_{IC}$ no debe usarse como una propiedad universal que cubra cualquier geometría, régimen y mecanismo.
\section{Qué información lleva realmente un factor geométrico}
La expresión $Y\sigma\sqrt{\pi a}$ separa una escala de tensión, una escala de tamaño y una corrección de geometría y carga. $Y$ puede depender de $a/W$, profundidad, relación de aspecto, forma de la superficie y distribución de tensión. Una grieta de borde y una grieta central no comparten automáticamente definición de $a$ ni factor.

Si $Y=Y(a)$, despejar $a_c$ con un valor fijo puede ser inconsistente. Se formula $f(a)=Y(a)\sigma\sqrt{\pi a}-K_c$ y se busca una raíz dentro del dominio físico. La comprobación posterior sustituye la raíz en la ecuación original y verifica ligamento y otras hipótesis. El problema P10 del cuaderno ilustra esta diferencia sin presentar su solución geométrica como válida para toda pieza.
\section{Del defecto observado a una decisión documentada}
La medida de un defecto incluye incertidumbre, orientación, acceso y método. Una indicación superficial puede requerir caracterización adicional; una longitud proyectada no determina siempre profundidad. El análisis usa una idealización explícita de defecto y justifica cómo se relaciona con la evidencia. Después considera las cargas máximas, el crecimiento entre inspecciones y la probabilidad de detección.

La tolerancia al daño integra resistencia residual, evolución del defecto y capacidad de inspección. Si una inspección no cubre una zona crítica, no se compensa solo reduciendo un tamaño inicial en la fórmula. La documentación aplicable determina la actuación real; el cálculo académico enseña a detectar qué hipótesis faltan antes de aceptar una conclusión.
\repaso{Si dos piezas tienen el mismo $K_I$, ¿tienen necesariamente igual seguridad?}{No. Pueden diferir tenacidad aplicable, espesor, ambiente, plasticidad, historia de carga, tensiones residuales y capacidad de detección. La igualdad del parámetro aplicado no iguala todos los mecanismos ni la evidencia.}

\resumen{\textbf{Preguntas de control.} $K_t$ es adimensional; $K_I$ tiene unidades de tensión por raíz de longitud. La longitud total de una grieta central es $2a$. $K_{IC}$ no es un esfuerzo aplicado. La comparación estática no excluye propagación subcrítica. El espesor y el ligamento forman parte de la validez.}
\fuentes{Roylance, \textit{Introduction to Fracture Mechanics}, apartados de Griffith, intensidad y tenacidad \cite{MIT-FRAC}; documentos docentes de fractura, pp. 1--46. El criterio dimensional se explica como cribado; no se afirma haber ejecutado un ensayo ASTM E399.}
